\chapter{Evaluation}\label{cha:evaluation}

This chapter presents the experimental results of the QBF Transformer tool. 
Section~\ref{sec:setup} describes the experimental setup, and Section~\ref{sec:equivalence-verification} verifies the correctness of the transformations via the equivalence rate. Sections~\ref{sec:runtime}--\ref{sec:secondary} present the measurements for runtime, coverage, and the relationship between coverage change and runtime effect. Section~\ref{sec:eval-summary} summarises all results in a single table.

\section{Experimental Setup}\label{sec:setup}

\subsection{Benchmark Formulas}\label{sec:benchmarks}

The experiments were conducted on a set of 50 QBF benchmark formulas in
QDIMACS format, drawn from the QBFLIB repository (Basler, Castellini,
Jordan-Kaiser, Kronegger-Pfandler-Pichler, Tacchella, Tentrup, and
Wintersteiger benchmark families).
Formulas were selected using the \texttt{filter\_benchmarks.py} utility,
retaining only instances that satisfy all of the following criteria:
\begin{itemize}
  \item at most 200 variables,
  \item at most 1\,000 clauses,
  \item at least two quantifier prefix blocks, and
  \item at least one universal and at least one existential quantifier block
        (formulas with only a single quantifier type are excluded as trivial).
\end{itemize}
 
A total of 400 experiments were recorded across all operator and formula
combinations (50 formulas $\times$ 8 operators).
After filtering for operator applicability, 285 experiments were retained
for analysis.
The equivalence rate and the reasons for exclusion are discussed in
Section~\ref{sec:equivalence-verification}.


\subsection{Operators and Variants}\label{sec:operators}

Eight operator variants were evaluated, grouped into four families:
 
\begin{itemize}
  \item \textbf{Tautology Insertion} (1 variant):
        \texttt{tautology\_insertion}
  \item \textbf{Pure Literal Generator} (3 variants, differing in
        which polarity of pure literal is selected):
        \texttt{pure\_literal\_gen\_auto} (any polarity),
        \texttt{pure\_literal\_gen\_pos} (positive only),
        \texttt{pure\_literal\_gen\_neg} (negative only)
  \item \textbf{Universal Reduction Trigger} (1 variant):
        \texttt{universal\_reduction\_trigger}
  \item \textbf{Blocked Clause Insertion} (3 variants, differing in the blocking-literal selection):
        \texttt{\seqsplit{blocked\_clause\_insertion}} (first existential, positive),
        \texttt{blocked\_clause\_insertion\_neg} (first existential, negated),
        \texttt{blocked\_clause\_insertion\_v1} (second existential)
\end{itemize}

\subsection{Measurement Parameters}\label{sec:measurement}

Each formula was solved once per solver per experiment. As the focus of this 
thesis is correctness verification and coverage analysis rather than fine-grained 
performance benchmarking, a single run is adequate to capture operator-specific 
effects. A per-experiment timeout of 60 seconds was applied. 
Runtime was measured with \texttt{time.perf\_counter}; code coverage was
collected via \texttt{gcov} on an instrumented build of DepQBF compiled
with \texttt{gcc-15}.


\subsection{Equivalence Verification}\label{sec:equivalence-verification}

Before analysing runtime and coverage, the correctness of each transformation
must be verified.
An experiment is marked as \emph{equivalent} if and only if DepQBF returns
the same result (SAT or UNSAT) for both the original and the mutant
formula (see Section~\ref{sec:equivalence-verification-impl}).
CAQE results were additionally recorded and showed no discrepancies when
inspected post-hoc, but are not used as part of the formal equivalence flag.

Of 400 total experiments, 285 (71.25\%) produced valid results and were
retained for runtime and coverage analysis.
The remaining 115 experiments were excluded for two reasons:


\begin{itemize}
  \item \textbf{Operator non-applicability} (99 experiments): the Pure
        Literal Generator variants could not be applied to formulas
        that contained no pure literals.
  \item \textbf{Parse errors} (16 experiments): two benchmark files
        (\texttt{p10-1.pddl\_planlen=0} and
        \texttt{p5-5.pddl\_planlen=0}) contained an empty quantifier
        block (\texttt{e~0}), which is non-standard QDIMACS and caused
        a parse error in the tool
        ($2~\text{files} \times 8~\text{operators} = 16~\text{experiments}$).
\end{itemize}

Crucially, all 285 retained experiments achieved a 100\% equivalence rate:
DepQBF returned identical results (SAT or UNSAT) for every
original--mutant pair, confirming that all four transformation operators
are correctly implemented and truth-value-preserving in practice.

Figure~\ref{fig:equivalence} shows the equivalence rate per operator.
Operators for which all experiments were retained (Tautology Insertion,
Universal Reduction Trigger, and all Blocked Clause variants) show a bar
at 100\% with $n = 48$.
The Pure Literal Generator variants show lower $n$ values (19, 13, 13)
reflecting the cases where the operator was not applicable.

\begin{figure}[h]
  \centering
  \includegraphics[width=\textwidth]{images/equivalence_rate.png}
  \caption{%
    Equivalence rate per transformation operator.
    The dashed green line marks the 100\% target.
    Bar labels show the equivalence rate and sample count~$n$.
    Operators with $n < 48$ were not applicable to all benchmark
    formulas (see text).%
  }
  \label{fig:equivalence}
\end{figure}

\section{Runtime Analysis}\label{sec:runtime}

The runtime effect of a transformation is quantified by the runtime ratio:
\[
  \text{ratio} = \frac{\text{median\_time}(\text{mutant})}{\text{median\_time}(\text{original})}.
\]
A ratio of $1.0$ indicates no change; values above $1.0$ indicate the mutant is slower; values below $1.0$ indicate the mutant is faster.

\subsection{Runtime Ratio per Operator — DepQBF}

Figure~\ref{fig:ratio-depqbf} shows the distribution of runtime ratios for DepQBF across all operators.

\begin{figure}[h]
  \centering
  \includegraphics[width=\textwidth]{images/time_ratio_boxplot_depqbf.png}
  \caption{Runtime ratio (mutant / original) per operator for DepQBF. The dashed red line marks ratio~$= 1.0$.}
  \label{fig:ratio-depqbf}
\end{figure}

\subsection{Runtime Ratio per Operator — CAQE}

Figure~\ref{fig:ratio-caqe} shows the same distribution for CAQE.

\begin{figure}[h]
  \centering
  \includegraphics[width=\textwidth]{images/time_ratio_boxplot_caqe.png}
  \caption{Runtime ratio (mutant / original) per operator for CAQE. The dashed red line marks ratio~$= 1.0$.}
  \label{fig:ratio-caqe}
\end{figure}

\subsection{Solver Comparison: DepQBF vs.\ CAQE}

Figure~\ref{fig:solver-comparison} directly compares the median runtime ratio of DepQBF and CAQE for each operator. Operators where both bars are on the same side of $1.0$ suggest a general structural effect; operators where the bars diverge suggest a solver-specific implementation effect.

\begin{figure}[h]
  \centering
  \includegraphics[width=\textwidth]{images/solver_comparison.png}
  \caption{Median runtime ratio comparison between DepQBF and CAQE per operator.}
  \label{fig:solver-comparison}
\end{figure}

\subsection{Interpretation}

The Tautology Insertion operator produces a median runtime ratio of 
$0.990$ for DepQBF and $0.999$ for CAQE, both close to~$1.0$. The values 
fall just below~$1.0$, consistent with measurement noise at this scale, 
and confirm that adding a trivially discarded clause causes only minimal overhead.
The Pure Literal Generator variants consistently show ratios below $1.0$ for DepQBF 
(auto: $0.968$, positive: $0.976$, negative: $0.980$), consistent with the prediction 
that eliminating clauses reduces the formula size and accelerates DepQBF. 
For CAQE the corresponding ratios are essentially at $1.0$ (auto: $1.002$, 
positive: $1.003$, negative: $0.999$), suggesting that the size reduction has 
a much weaker effect on CAQE's clausal-abstraction strategy than on DepQBF's 
search-based procedure. The Universal Reduction Trigger shows a small but 
notable solver asymmetry: $0.988$ for DepQBF (mutant slightly faster) 
versus $1.005$ for CAQE (mutant slightly slower), suggesting that 
DepQBF benefits marginally from the inserted clause being simplified 
away during preprocessing, while CAQE incurs a small overhead from 
the additional clause. The Blocked Clause Insertion variants show 
ratios slightly below $1.0$ for DepQBF ($0.988$--$0.992$), and ratios 
slightly above $1.0$ for CAQE ($1.002$--$1.005$), consistent with 
the same solver asymmetry and the formula-structure dependence of 
blocked clause detection.

\section{Coverage Analysis}\label{sec:coverage}

Code coverage is measured as the percentage of DepQBF source lines and branches
executed during solving. The coverage delta~$\Delta$ is defined as:
\[
  \Delta = \text{coverage}(\text{mutant}) - \text{coverage}(\text{original}).
\]
A positive delta indicates that the mutant formula caused the solver to execute
additional code; a negative delta indicates that fewer code paths were reached.
Coverage was collected only for DepQBF, as \texttt{gcov} instrumentation is not
available for the CAQE binary.

\subsection{Line Coverage Delta}

Figure~\ref{fig:coverage-line} shows the median line coverage delta per operator.
Error bars represent the interquartile range~(IQR) across all experiments in that
operator group.

\begin{figure}[h]
  \centering
  \includegraphics[width=\textwidth]{images/coverage_delta_line.png}
  \caption{%
    Line coverage delta ($\Delta$ line coverage in percentage points,~pp) per
    operator. Bars show the median; error bars show the interquartile range~(IQR).
    Positive values indicate the mutant activates additional source lines in DepQBF.%
    }
  \label{fig:coverage-line}
\end{figure}

\subsection{Branch Coverage Delta}

Figure~\ref{fig:coverage-branch} shows the median branch coverage delta per operator.
Branch coverage is strictly finer than line coverage: a single \texttt{if}-statement
contributes two branches, so branch deltas tend to be slightly larger in magnitude
than the corresponding line deltas.

\begin{figure}[h]
  \centering
  \includegraphics[width=\textwidth]{images/coverage_delta_branch.png}
  \caption{%
    Branch coverage delta ($\Delta$ branch coverage in percentage points,~pp) per
    operator. Bars show the median; error bars show the interquartile range~(IQR).
    Positive values indicate the mutant causes additional conditional branches to be
    taken in DepQBF.%
     }
  \label{fig:coverage-branch}
\end{figure}

\subsection{Interpretation}\label{sec:coverage-interpretation}

\textbf{Tautology Insertion} produces a median line coverage delta of
$+0.033\,\text{pp}$ and a branch delta of $+0.030\,\text{pp}$.
Both values are positive, indicating that the inserted tautological clause
$(\ell \lor \neg\ell)$ triggers an additional preprocessing code path in DepQBF
that is not reached by the original formula.
 
\textbf{Pure Literal Generator} (all three variants: auto, positive, negative)
consistently produce negative deltas of $-0.017\,\text{pp}$ (line) and
$-0.030\,\text{pp}$ (branch).
This is consistent with the design intent of this operator: by applying Pure Literal
Elimination to produce a strictly smaller mutant, the transformed formula traverses
fewer solver-internal code paths than the original. All three variants produce
identical coverage deltas because they share the same elimination mechanism and
differ only in which pure literal is selected.

\textbf{Universal Reduction Trigger} produces a median line delta of
$+0.033\,\text{pp}$ and branch delta of $+0.030\,\text{pp}$, identical to
Tautology Insertion.
This suggests that the inserted clause $[u, e, \neg e]$ activates the same
preprocessing code path as the tautological clause — consistent with the
expectation that Universal Reduction is applied during the same preprocessing
phase as tautology detection.
 
\textbf{Blocked Clause Insertion} (Standard variant) produces a median line
delta of $+0.017\,\text{pp}$ and branch delta of $+0.015\,\text{pp}$,
approximately half the delta of Tautology Insertion. This suggests partial
activation of the blocked clause detection code path.
 
\textbf{Blocked Clause (Negated)} produces the largest positive delta of all
operators: $+0.075\,\text{pp}$ line and $+0.091\,\text{pp}$ branch. The larger
delta compared to the standard variant suggests that negating the blocking literal
activates a different, less frequently reached branch within the blocked clause
detection routine of DepQBF.
 
\textbf{Blocked Clause (Variant~2)} produces a delta of exactly $0\,\text{pp}$
for both line and branch coverage. The inserted clause is handled by an
already-covered code path, meaning no new lines or branches are reached by the mutant.

It is noteworthy that the Blocked Clause (Negated) and Blocked Clause (Standard)
variants exhibit large IQR error bars (approximately $\pm 0.5\,\text{pp}$),
reflecting high variability across benchmark formulas. This variability is
formula-structure-dependent: blocked clause detection is only triggered when the
inserted clause is blocked with respect to the formula's existing clause set,
which varies per instance.

\section{Secondary Analysis: Coverage Change vs.\ Runtime Effect}\label{sec:secondary}

This section investigates whether a larger coverage delta is associated with a larger runtime ratio. Figures~\ref{fig:cov-vs-ratio-depqbf} and~\ref{fig:cov-vs-ratio-caqe} provide an overview across both solvers: each point represents one experiment, with the coverage delta on the $x$-axis and the runtime ratio on the $y$-axis. Points are colour-coded by operator family.

\begin{figure}[h]
  \centering
  \includegraphics[width=\textwidth]{images/scatter_depqbf.png}
  \caption{Coverage delta (pp) vs.\ runtime ratio for DepQBF. The dashed red line marks ratio $= 1.0$.}
  \label{fig:cov-vs-ratio-depqbf}
\end{figure}

\begin{figure}[h]
  \centering
  \includegraphics[width=\textwidth]{images/scatter_caqe.png}
  \caption{Coverage delta (pp) vs.\ runtime ratio for CAQE. The dashed red line marks ratio $= 1.0$.}
  \label{fig:cov-vs-ratio-caqe}
\end{figure}

Figure~\ref{fig:coverage-vs-ratio} then breaks the DepQBF data down by operator family, with a linear trend line with slope $m$ fitted per group.

\begin{figure}[h]
  \centering
  \includegraphics[width=\textwidth]{images/coverage_vs_ratio_depqbf.png}
  \caption{Coverage delta (pp) vs.\ runtime ratio for DepQBF, grouped by operator family. The dashed coloured line shows a linear fit with slope $m$; the dotted red line marks ratio $= 1.0$.}
  \label{fig:coverage-vs-ratio}
\end{figure}

For three of the four operator groups the fitted slope is close to zero
($|m| \leq 0.074$), indicating no meaningful linear relationship between
coverage delta and runtime ratio.
The exception is \textbf{Tautology Insertion} ($m = -0.472$): a larger
positive coverage delta is weakly associated with a \emph{lower} runtime
ratio, suggesting that reaching additional preprocessing code paths
slightly reduces solving time rather than increasing it.
However, the Tautology Insertion group spans a narrow coverage range
(approximately $-0.01$ to $+0.03\,\text{pp}$) with few distinct delta
values, so this slope should not be over-interpreted.

The \textbf{Blocked Clause Addition} group exhibits the widest spread on
both axes: coverage deltas range from approximately $-1\,\text{pp}$ to
$+3\,\text{pp}$, and runtime ratios range from roughly $0.75$ to $1.25$.
Despite this spread, the slope is nearly flat ($m = 0.003$), confirming
that coverage delta is not a predictor of runtime change within this group.
The large variance is consistent with the formula-structure dependence of
blocked clause detection noted in Section~\ref{sec:coverage-interpretation}.

Overall, the analysis provides no evidence that a larger coverage delta
causes a larger runtime effect across any operator group.

\section{Summary}\label{sec:eval-summary}

\begin{table}[h]
\centering
\caption{Aggregated results per transformation operator (n = number of valid experiments; ratios and coverage deltas are medians across formulas).}
\label{tab:summary}
\resizebox{\textwidth}{!}{%
\begin{tabular}{llrrrrr}
\toprule
\textbf{Operator} & \textbf{Group} & \textbf{n} &
\textbf{Ratio DepQBF} & \textbf{Ratio CAQE} &
\textbf{$\Delta$ Line (pp)} & \textbf{$\Delta$ Branch (pp)} \\
\midrule
\texttt{tautology\_insertion}            & Tautology           & 48 & $0.990$ & $0.999$ & $+0.033$ & $+0.030$ \\
\texttt{pure\_literal\_gen\_auto}        & Pure Literal        & 19 & $0.968$ & $1.002$ & $-0.017$ & $-0.030$ \\
\texttt{pure\_literal\_gen\_pos}         & Pure Literal        & 13 & $0.976$ & $1.003$ & $-0.017$ & $-0.030$ \\
\texttt{pure\_literal\_gen\_neg}         & Pure Literal        & 13 & $0.980$ & $0.999$ & $-0.017$ & $-0.030$ \\
\texttt{universal\_reduction\_trigger}   & Universal Reduction & 48 & $0.988$ & $1.005$ & $+0.033$ & $+0.030$ \\
\texttt{blocked\_clause\_insertion}      & Blocked Clause      & 48 & $0.988$ & $1.005$ & $+0.017$ & $+0.015$ \\
\texttt{blocked\_clause\_insertion\_neg} & Blocked Clause      & 48 & $0.992$ & $1.002$ & $+0.075$ & $+0.091$ \\
\texttt{blocked\_clause\_insertion\_v1}  & Blocked Clause      & 48 & $0.989$ & $1.004$ & $+0.000$ & $+0.000$ \\
\bottomrule
\end{tabular}}
\end{table}

All DepQBF runtime ratios fall within the range $0.968$--$0.992$, confirming
that none of the transformations causes measurable overhead for DepQBF\@.
The Pure Literal Generator variants produce the lowest DepQBF ratios
($0.968$--$0.980$), consistent with the expectation that a smaller mutant
formula reduces solving time.

For CAQE, the picture is reversed for several operators: Universal Reduction,
Blocked Clause (Standard), and Blocked Clause (Variant~2) all show CAQE ratios
slightly above $1.0$ ($1.004$--$1.005$), meaning the mutant is marginally
\emph{slower} for CAQE\@. This solver-specific asymmetry suggests that the two
solvers handle the inserted clauses differently during preprocessing: DepQBF
appears to simplify them away at a net benefit, while CAQE incurs a small
overhead.

Coverage deltas remain small across all operators (below $0.10\,\text{pp}$
in all cases), reflecting the limited size of the benchmark formulas.
The Blocked Clause (Negated) variant produces the largest positive delta
($+0.075\,\text{pp}$ line, $+0.091\,\text{pp}$ branch), suggesting that
negating the blocking literal activates additional preprocessing code paths
in DepQBF\@. The Pure Literal Generator variants consistently produce negative
deltas ($-0.017\,\text{pp}$ line, $-0.030\,\text{pp}$ branch); the three
variants yield identical deltas because they share the same elimination
mechanism and differ only in which pure literal is selected. A smaller
mutant formula causes the solver to traverse fewer code paths overall.